\documentclass[12pt]{article}
\usepackage[utf8]{inputenc}
\usepackage[T1]{fontenc}
\usepackage{hyperref}
\title{So'a Language Resources for K-12 Curriculum Development}
\author{Compiled by Claude Sonnet 5.5 (claude-sonnet-5-5) for Andr\'as Kornai}
\date{}
\begin{document}
\maketitle

\section{Introduction}
So'a (also Soa; sometimes called Central Ngada; ISO 639-3 ssq) is an Austronesian language of central Flores, East Nusa Tenggara, Indonesia. It forms a dialect cluster with Ngadha and lies between the Ngad'a and Riung language areas. Wikipedia gives about 10,000 native speakers (a 1994 figure). Ethnologue rates it vulnerable, with low confidence. Very little material was found, and nothing indicates that So'a is taught in schools.

\section{Primary Educational Resources}
No primers, textbooks or school programmes for So'a were found.

\section{Linguistic and Reference Materials}
\begin{itemize}
\item Overview of the language: \url{https://en.wikipedia.org/wiki/So%CA%BCa_language}
\item Endangered Languages Project entry: \url{https://endangeredlanguages.com/lang/9589?hl=en}
\item Verbix pages: \url{https://docs.verbix.com/Languages/Soa}
\end{itemize}
No grammar or dictionary was identified in the search results.

\section{Religious and Community Resources}
Joshua Project notes that the So'a people are mostly Christian, Catholicism being the most common tradition: \url{https://joshuaproject.net/people_groups/14968}. No scripture translation or community organisation was found.

\section{Digital and Audio Resources}
None identified.

\section{Curriculum Development Context}
Documentation is minimal. Any native-language curriculum would need new documentation work with the community, an agreed orthography, and coordination with the closely related Ngadha. Indonesian is the likely language of schooling; no source on local-language policy was found in these searches.

\end{document}